\documentclass[10pt,a4paper]{article}
\usepackage[margin=18mm]{geometry}
\usepackage{booktabs,amsmath,array}
\setlength{\parindent}{0pt}
\setlength{\parskip}{7pt}
\begin{document}
\begin{center}
{\large\bfseries Repository Supporting Tables}\\[5pt]
Reliability-First Evaluation for Spatio-Temporal Graph Forecasting
\end{center}
\section*{Table S1. Complete reference-degree controls}
The ten cases below report changes in flat-pooled target MAE after zeroing the
selected ten sensors across all input channels and history steps. Each case uses
its own clean epoch-100 checkpoint reference. Uniform and degree-control values
are means and sample SDs across 30 random draws, not retraining variability.
The selected sensor set is fixed. The rank is one plus the number of degree-control
draws with a smaller MAE change, among the selected value and 30 controls.
Higher ranks indicate larger selected-set MAE changes relative to the degree
controls. Rank 31 means a larger increase than all 30 controls, not better forecast
accuracy. Ranks are descriptive, not significance tests.

\begin{center}\small
\setlength{\tabcolsep}{3pt}
\begin{tabular}{llrrrrr}
\toprule
Dataset & Model & Clean MAE & Selected $\Delta$ & Uniform $\Delta\pm$SD & Degree $\Delta\pm$SD & Rank\\
\midrule
METR-LA & D2STGNN & 2.879019 & 0.307359 & $0.302450\pm0.073074$ & $0.291343\pm0.069949$ & 22 \\
PEMS-BAY & D2STGNN & 1.514080 & 0.100838 & $0.124713\pm0.015541$ & $0.125062\pm0.018997$ & 5 \\
PEMS04 & D2STGNN & 18.402961 & 4.105417 & $3.334588\pm0.598201$ & $3.178515\pm0.393061$ & 31 \\
METR-LA & MTGNN & 3.065592 & 0.670083 & $0.298490\pm0.032341$ & $0.338463\pm0.056788$ & 31 \\
PEMS-BAY & MTGNN & 1.602254 & 0.140866 & $0.115981\pm0.014924$ & $0.114733\pm0.016701$ & 27 \\
PEMS04 & MTGNN & 19.098077 & 2.190102 & $2.147260\pm0.328576$ & $2.174683\pm0.343091$ & 17 \\
METR-LA & MegaCRN & 3.149723 & 0.176190 & $0.298734\pm0.102925$ & $0.318347\pm0.114101$ & 4 \\
METR-LA & STGCN-Cheb & 3.162829 & 0.071302 & $0.133999\pm0.032916$ & $0.126184\pm0.029872$ & 1 \\
PEMS-BAY & STID & 1.563499 & 0.146503 & $0.125829\pm0.016737$ & $0.114087\pm0.014972$ & 30 \\
PEMS04 & STID & 18.427375 & 2.072549 & $3.471690\pm0.543054$ & $3.685005\pm0.561158$ & 1 \\
\bottomrule
\end{tabular}
\end{center}
Degree is measured in the supplied reference topology, not each architecture's
learned or dynamic graph. Nine cases match selected degrees exactly.
MTGNN on METR-LA uses degree 15 for one selected degree-18 slot because no
unselected equal-degree sensor is available. The mean absolute degree mismatch
is 0.3 across the ten slots.

These controls characterize checkpoint-specific target-error sensitivity and are
interpreted separately from the sample-specific forecast-change diagnostics in
the main manuscript.
\texttt{Table S1.csv} identifies each source JSON, its SHA-256 and the rank field
in the unchanged technical release \texttt{v1.4.4}.

\clearpage
\section*{Table S2. Calibration-boundary validity-mask check}
This saved-array calculation keeps the original 3,415 calibration origins,
ensemble members, fitted thresholds, interval definitions and target-validity
mask unchanged. It excludes the first 11 evaluation origins, reducing the
evaluation from 3,416 to 3,405 origins. No training or model inference is performed.

The first 11 evaluation origins contribute no valid target elements under the
predefined mask because their target values are zero in the supplied arrays.
The mask retains finite targets unequal to zero. These values match the
corresponding supplied benchmark-array windows. Both evaluations therefore
contain the same 7,118,436 valid target elements. Exclusion changes the origin
count, not the evaluated target set.

\begin{center}\small
\begin{tabular}{lrrrr}
\toprule
 & \multicolumn{2}{c}{PICP (\%)} & \multicolumn{2}{c}{MPIW}\\
Method & Original & Excluded & Original & Excluded\\
\midrule
Dispersion-normalized, pooled & 90.56704 & 90.56704 & 23.378166 & 23.378166 \\
Dispersion-normalized, per horizon & 90.56291 & 90.56291 & 23.180659 & 23.180659 \\
Unscaled, pooled & 90.14365 & 90.14365 & 13.789711 & 13.789711 \\
\bottomrule
\end{tabular}
\end{center}
MPIW uses the dataset's target units. Aggregate and all 12 horizon-wise PICP and
MPIW values are unchanged for all three methods. The complete values are in
\texttt{Table S2 all horizons.csv}. The corresponding JSON records the array
hashes, metric comparisons and calibration quantiles. Recalculation reproduces
the recorded aggregate and horizon-wise metrics.

This checks boundary handling for the retained METR-LA arrays and mask, not
sensitivity to removing overlapping valid observations. It does not establish
independence, exchangeability or conditional-coverage guarantees, or test the
other datasets' boundary windows.
\end{document}
